\section{Motivation}\label{sec:motivation}

Respiratory motion is a well-known complicating factor in MRI, especially when imaging elements of the cardiovascular system. Periodic motion of the lungs during the respiratory cycle affects the torso and abdomen, which impacts imaging of the heart, great vessels, and the liver. The resulting motion artifacts, which appear in traditional Cartesian imaging as distinct replicas in the phase-encoding direction, negatively impact the quality of diagnostic information that can be gleaned from MR images. There are many approaches to mitigate the effects of respiratory motion, including prospectively gating the acquisition to only acquire during a specific phase of the respiratory cycle or asking the patient to hold their breath for the duration of the acquisition1. However, respiration is also a natural part of human physiology and there is valuable information that can be extracted from the full respiratory cycle.

\section{Hypothesis}\label{sec:hypothesis}
In this thesis, I seek to implement and optimize methods for not only mitigating the effects of respiration in cardiovascular MRI which would enable otherwise impossible imaging, but to also utilize information from the respiratory cycle to improve image quality, increase data utilization efficiency, and extract valuable diagnostic information. This will be performed in three different scientific contexts, each focusing on a different vascular region: 1.\ improving and optimizing radial free-breathing 2D Phase Contrast MRI for measuring aortic pulse wave velocity, 2.\ accelerating the acquisition of radial 4D flow MRI in the portal venous system, and 3.\ implementing a 5D MRI approach for resolving cardiac and respiratory motion of the heart in radiation therapy.

\section{Chapter Summary}
\begin{description}
    \item \textbf{Chapter 2}
    \item \textbf{Chapter 3}
    \item \textbf{Chapter 4}
    \item \textbf{Chapter 5}
    \item \textbf{Chapter 6}
    \item \textbf{Chapter 7}
    \item \textbf{Chapter 8}
\end{description}

